\documentclass[11pt,a4paper]{article}

% -------------------------------------------------
% Encoding / Language
% -------------------------------------------------
\usepackage[T1]{fontenc}
\usepackage[utf8]{inputenc}
\usepackage[english]{babel}

% -------------------------------------------------
% Math
% -------------------------------------------------
\usepackage{amsmath,amssymb,amsthm,mathtools}
\usepackage{bm}

% -------------------------------------------------
% Layout
% -------------------------------------------------
\usepackage[a4paper,margin=2.5cm]{geometry}
\usepackage{hyperref}

% -------------------------------------------------
% Theorems
% -------------------------------------------------
\newtheorem{definition}{Definition}[section]
\newtheorem{theorem}[definition]{Theorem}
\newtheorem{lemma}[definition]{Lemma}
\newtheorem{remark}[definition]{Remark}

% -------------------------------------------------
% Title
% -------------------------------------------------
\title{
	Global Stability vs.\ Directional Positivity in the 2HDM \\
	\large A Structural Correction and a Coupling-Based Framework
}

\author{
	Thomas Hoffbauer \\
	\small Bamberg, Germany
}

\date{}

\begin{document}
	
	\maketitle
	
	\begin{abstract}
		A recent formal verification of the two Higgs doublet model (2HDM) has revealed an error in a widely cited stability criterion. The error originates from replacing a global boundedness condition with local directional constraints. In this paper, we show that the correct condition follows directly from quadratic completion and necessarily involves a coupling inequality between linear and quadratic contributions. We further propose a structural interpretation in which stability is understood as a resonance condition in a coupled system. This perspective naturally extends to an operator-theoretic framework.
	\end{abstract}
	
	\section{Introduction}
	
	The stability of scalar potentials is a central problem in quantum field theory. In the context of the 2HDM, a widely used condition for stability has recently been shown to be incorrect via formal verification in Lean.
	
	The issue is conceptual: a global condition has been replaced by local directional checks. This paper argues that the mistake is structurally visible and that the correct condition follows from elementary analysis.
	
	\section{Reduction to a One-Parameter Problem}
	
	The 2HDM potential can be written as a function of a scaling parameter $K_0 > 0$ and a normalized direction $\vec{k}$ with $\|\vec{k}\|^2 \leq 1$:
	\begin{equation}
		V(K_0,\vec{k}) = K_0 J_2(\vec{k}) + K_0^2 J_4(\vec{k}).
	\end{equation}
	
	Thus, for fixed $\vec{k}$, stability reduces to a one-dimensional problem.
	
	\section{Stability as a Global Property}
	
	\begin{definition}
		The potential is stable if there exists $c \in \mathbb{R}$ such that
		\[
		c \leq V(K_0,\vec{k})
		\quad \text{for all } K_0 > 0, \ \|\vec{k}\|^2 \leq 1.
		\]
	\end{definition}
	
	This involves a mixed quantifier structure:
	\[
	\exists c \ \forall K_0 \ \forall \vec{k}.
	\]
	
	\section{Quadratic Completion}
	
	\begin{lemma}
		For fixed $\vec{k}$, the function
		\[
		f(K_0) = K_0 J_2(\vec{k}) + K_0^2 J_4(\vec{k})
		\]
		is bounded from below if and only if:
		\[
		J_4(\vec{k}) \geq 0
		\quad \text{and} \quad
		J_2(\vec{k})^2 \leq 4c\,J_4(\vec{k})
		\]
		for some $c \geq 0$.
	\end{lemma}
	
	\begin{proof}
		Completing the square:
		\[
		f(K_0) = J_4(\vec{k})\left(K_0 + \frac{J_2(\vec{k})}{2J_4(\vec{k})}\right)^2 
		- \frac{J_2(\vec{k})^2}{4J_4(\vec{k})}.
		\]
		The minimum is finite if and only if $J_4(\vec{k}) \geq 0$, and the lower bound is controlled by the second term.
	\end{proof}
	
	\section{Failure of Directional Positivity}
	
	The incorrect claim replaces global boundedness with purely directional conditions on $J_2$ and $J_4$.
	
	However, the dependence on $K_0$ cannot be eliminated.
	
	\begin{theorem}
		Directional positivity conditions are not sufficient for global stability.
	\end{theorem}
	
	\begin{remark}
		This has been formally demonstrated by constructing explicit counterexamples satisfying directional constraints but violating boundedness. :contentReference[oaicite:0]{index=0}
	\end{remark}
	
	\section{Coupling Interpretation}
	
	The potential naturally decomposes into:
	\begin{itemize}
		\item Linear term: $J_2(\vec{k})$
		\item Quadratic term: $J_4(\vec{k})$
	\end{itemize}
	
	Stability requires not only positivity but a coupling condition:
	\[
	J_2^2 \lesssim J_4.
	\]
	
	\section{Bamberg Model Perspective}
	
	We reinterpret the structure dynamically:
	
	\begin{itemize}
		\item $J_2$: drift (Tri-component)
		\item $J_4$: curvature (Okto-component)
		\item $K_0$: scaling (energy axis)
	\end{itemize}
	
	\begin{definition}
		A system is stable if its components satisfy a resonance condition:
		\[
		(\text{Tri})^2 \lesssim (\text{Okto}).
		\]
	\end{definition}
	
	Thus, stability is not a geometric property of directions but a dynamical balance across scales.
	
	\section{Operator-Theoretic Outlook}
	
	The structure suggests an operator formulation:
	
	\begin{equation}
		H = K_0 D_1 + K_0^2 D_2,
	\end{equation}
	
	where $D_1$ and $D_2$ encode linear and quadratic contributions.
	
	Stability corresponds to spectral boundedness of $H$, which requires compatibility between $D_1$ and $D_2$.
	
	This aligns with Hilbert--Pólya-type ideas, where spectral properties encode global constraints.
	
	\section{Conclusion}
	
	The error identified through formalization reflects a structural issue: replacing global constraints with local checks.
	
	We have shown that:
	
	\begin{itemize}
		\item Stability reduces to a quadratic minimization problem
		\item The correct condition is a coupling inequality
		\item Stability is inherently a resonance phenomenon
	\end{itemize}
	
	\bigskip
	
	\noindent
	\textbf{Key Insight:} Stability is not local positivity, but global coupling.
	
\end{document}